\documentclass[11pt,letterpaper,oneside]{report}
\usepackage[T1]{fontenc}
\usepackage[utf8]{inputenc}
\usepackage{lmodern}
\usepackage[letterpaper,margin=0.92in,headheight=15pt]{geometry}
\usepackage{amsmath,amssymb,amsthm,mathtools}
\usepackage{microtype}
\usepackage[dvipsnames]{xcolor}
\usepackage{graphicx}
\usepackage{booktabs,longtable,array,tabularx}
\usepackage{enumitem}
\usepackage{fancyhdr}
\usepackage{titlesec}
\usepackage[most]{tcolorbox}
\usepackage{listings}
\usepackage{tikz}
\usetikzlibrary{positioning}
\usepackage{hyperref}
\hypersetup{colorlinks=true,linkcolor=MidnightBlue,urlcolor=RoyalBlue,pdfauthor={Independent Mathematics and Engineering Study Suite},pdftitle={Mathematical Foundations and Prerequisites}}
\definecolor{ink}{HTML}{162639}
\definecolor{navy}{HTML}{173B61}
\definecolor{blue}{HTML}{2874A6}
\definecolor{teal}{HTML}{007C83}
\definecolor{amber}{HTML}{8B5E16}
\definecolor{pale}{HTML}{EDF4F8}
\definecolor{paleteal}{HTML}{ECF8F5}
\definecolor{paleamber}{HTML}{FFF6E8}
\titleformat{\chapter}[display]{\normalfont\huge\bfseries\color{navy}}{\filleft\fontsize{43}{46}\selectfont\thechapter}{0.4ex}{\titlerule\vspace{1ex}}[\vspace{1ex}\titlerule]
\titleformat{\section}{\normalfont\Large\bfseries\color{navy}}{\thesection}{0.8em}{}
\titleformat{\subsection}{\normalfont\large\bfseries\color{teal}}{\thesubsection}{0.75em}{}
\pagestyle{fancy}\fancyhf{}\fancyhead[L]{\small\color{navy}Differential Equations Engineering Suite}\fancyhead[R]{\small\color{navy}Volume 1 of 8}\fancyfoot[C]{\thepage}\renewcommand{\headrulewidth}{0.3pt}
\setlist[itemize]{leftmargin=*,itemsep=3pt,topsep=5pt}\setlist[enumerate]{leftmargin=*,itemsep=4pt,topsep=5pt}
\setlength{\parskip}{5pt}\setlength{\parindent}{0pt}\setlength{\emergencystretch}{2em}
\newtheorem{theorem}{Theorem}[chapter]
\newtheorem{proposition}[theorem]{Proposition}
\newtheorem{lemma}[theorem]{Lemma}
\theoremstyle{definition}\newtheorem{definition}[theorem]{Definition}
\newtheorem{example}[theorem]{Worked example}
\theoremstyle{remark}\newtheorem{remark}[theorem]{Remark}
\newtcolorbox{bookground}[1][]{enhanced,breakable,colback=pale,colframe=blue,boxrule=0.6pt,title={TEXTBOOK ALIGNMENT},fonttitle=\sffamily\bfseries\small,#1}
\newtcolorbox{modern}[1][]{enhanced,breakable,colback=paleteal,colframe=teal,boxrule=0.6pt,title={MODERN COMPUTATIONAL EXTENSION},fonttitle=\sffamily\bfseries\small,#1}
\newtcolorbox{caution}[1][]{enhanced,breakable,colback=paleamber,colframe=amber,boxrule=0.6pt,title={ASSUMPTION / VALIDATION},fonttitle=\sffamily\bfseries\small,#1}
\newcommand{\dd}{\,\mathrm{d}}
\newcommand{\R}{\mathbb{R}}
\newcommand{\lap}{\mathcal{L}}
\newcommand{\norm}[1]{\left\lVert#1\right\rVert}
\newcommand{\abs}[1]{\left\lvert#1\right\rvert}
\lstdefinestyle{engpy}{language=Python,basicstyle=\ttfamily\footnotesize,keywordstyle=\color{RoyalBlue}\bfseries,commentstyle=\color{gray},stringstyle=\color{teal},backgroundcolor=\color{pale},frame=single,rulecolor=\color{blue!30},breaklines=true,columns=fullflexible,keepspaces=true,showstringspaces=false,tabsize=4,upquote=true}
\lstset{style=engpy}
\begin{document}
\begin{titlepage}
\pagecolor{navy}\color{white}
\vspace*{0.75in}
{\sffamily\bfseries\large MATHEMATICS \& ENGINEERING DOCUMENTATION SUITE\\[12pt]}
{\color{cyan!40}\rule{\textwidth}{2pt}}\\[1.0in]
{\sffamily\bfseries\fontsize{29}{36}\selectfont Mathematical Foundations and Prerequisites\par}
\vspace{0.45in}
{\Large Volume 1 of 8}\par\vspace{0.55in}
{\color{cyan!30}\rule{0.50\textwidth}{1pt}}\par\vspace{0.3in}
{\large Independent technical companion aligned with the\par}
{\large 13-chapter organization of Nagle, Saff \& Snider (7th ed.)\par}
\vfill
{\large Original derivations, worked examples, verification and engineering practice\par}\vspace{0.3in}
{\small Prepared October 2026 \quad\textbullet\quad pdf\LaTeX{} compatible\par}
\end{titlepage}
\nopagecolor\color{ink}\pagenumbering{roman}
\chapter*{Reader's guide and source boundaries}\addcontentsline{toc}{chapter}{Reader's guide and source boundaries}
This independent study resource is organized around \emph{Fundamentals of Differential Equations and Boundary Value Problems}, Seventh Edition, by R.~Kent Nagle, Edward~B. Saff, and Arthur~David Snider (Pearson, 2018; ISBN 978-0-321-97710-6). The book's chapter and section numbers are given as \emph{navigation references}, not as claims of reproduced text or proofs.

\begin{bookground}
\textbf{Textbook-grounded scope:} topics, terminology, and chapter/section organization corresponding to the cited textbook. All explanations, derivations, exercise prompts, and diagrams in this volume are newly authored and may emphasize different examples. Check the textbook itself for its exact hypotheses, proofs, exercises, or nomenclature in a particular section.
\end{bookground}
\begin{modern}
\textbf{Distinct extension layer:} Python/NumPy/SciPy workflows, engineering QA and reproducibility practices, automated convergence studies, diagnostic plots, and implementation guidance. These are \emph{not} represented as content supplied by the textbook. API calls are examples and should be verified against the installed package version.
\end{modern}
\begin{caution}
\textbf{Mathematical qualification:} formulas require their stated regularity, domain and boundary/initial conditions. Numerical agreement does not replace a proof; computational outputs are not certified error bounds unless explicitly established.
\end{caution}
\tableofcontents\clearpage\pagenumbering{arabic}
\chapter{Differential equations as mathematical models}
\begin{bookground}Correspondence: Chapter 1 (Introduction), especially Sections 1.1--1.4; related modeling in Chapter 3 and existence questions in Chapter 13.\end{bookground}
\section{Objects, notation, and classification}
An \emph{ordinary differential equation} (ODE) constrains an unknown function of one independent variable and its derivatives. A \emph{partial differential equation} (PDE) constrains partial derivatives of a function of several variables. The order is the highest derivative appearing. A linear ODE has the form
\begin{equation}\label{eq:linear-n}a_n(x)y^{(n)}+\cdots+a_1(x)y'+a_0(x)y=g(x),\qquad a_n(x)\ne 0.\end{equation}
Its coefficients may depend on the independent variable but not on $y$ or its derivatives. A \emph{homogeneous} linear equation has $g\equiv0$; an equation with $g\not\equiv0$ is nonhomogeneous. \emph{Autonomous} systems $\mathbf x'=\mathbf f(\mathbf x)$ do not explicitly depend on time. A nonlinear differential equation can still have a linearization near an equilibrium.

\begin{example}[Classification and dependent variables]
The equation $x^2y''+xy'+4y=0$ is a second-order, homogeneous, variable-coefficient linear ODE for $y(x)$ on any interval excluding $x=0$. The equation $y'=y(1-y/K)$ is nonlinear, first order and autonomous. The equation $u_t=\kappa u_{xx}$ is a linear evolution PDE; the variables $t$ and $x$ are independent.
\end{example}
\section{Initial-value and boundary-value problems}
An \emph{initial-value problem} (IVP) combines an equation with initial data, e.g. $y'=f(t,y)$ and $y(t_0)=y_0$. A \emph{two-point boundary-value problem} (BVP) prescribes conditions at both ends of a spatial interval, e.g.
\[y''+\lambda y=0,\quad 0<x<L,\qquad y(0)=y(L)=0.\]
The distinction matters: an IVP is propagated from known state; a BVP is solved globally under endpoint constraints. Existence, uniqueness, and numerical conditioning can differ fundamentally.
\section{Solution, implicit solution, and interval of validity}
A solution must possess the derivatives required by the equation and satisfy it on a specified interval. A \emph{general solution} describes a family with arbitrary constants; a \emph{particular solution} selects data. An \emph{implicit solution} such as $y^2+x^2=C$ may describe branches requiring additional care. Algebraic manipulation can lose constant or singular solutions, and dividing by a function may exclude zeros. Always record the domain where coefficients and transformations are defined.
\begin{example}[Interval-limited solution]
For $y'=y^2$ and $y(0)=1$, separating gives $-1/y=t+C$, hence $y(t)=1/(1-t)$. This satisfies $y'=1/(1-t)^2=y^2$, but its maximal interval containing $0$ is $(-\infty,1)$. A numerical integrator cannot continue the finite solution past $t=1$.
\end{example}
\section{Direction fields and qualitative information}
For $y'=f(t,y)$, each point $(t,y)$ receives a short segment with slope $f(t,y)$. An \emph{isocline} has constant slope, and an autonomous scalar equilibrium satisfies $f(y_*)=0$. Without solving $y'=y(1-y)$, the sign chart yields increasing motion for $0<y<1$ and decreasing motion for $y>1$; $y=1$ is locally attracting, while $y=0$ repels positive solutions.
\chapter{Calculus prerequisite review}
\begin{bookground}Preparation supporting Chapters 1--4, 7--8, and Appendix A (Review of Integration Techniques). The organization here is an original prerequisite map.\end{bookground}
\section{Differentiation and the chain rule}
The chain rule $\frac{d}{dt}F(g(t))=F'(g(t))g'(t)$ underlies substitutions, autonomous equations and composition of solutions. The product rule gives $(\mu y)'=\mu y'+\mu' y$, which is the mechanism behind integrating factors. For $f(x,y(x))$, total differentiation yields
\[\frac{d}{dx}f(x,y(x))=f_x(x,y)+f_y(x,y)y'.\]
For multiple state variables $\mathbf x(t)$, the Jacobian $J(\mathbf x)$ is the first derivative of a vector field; near an equilibrium $\mathbf x_*+\mathbf u$,
\[\mathbf u'=J(\mathbf x_*)\mathbf u+O(\norm{\mathbf u}^2),\]
provided adequate smoothness.
\section{Integration techniques}
Separation uses $\int\!dy/g(y)=\int\!f(x)\dd x$ where $g(y)\ne0$. Integration by parts, $\int u\dd v=uv-\int v\dd u$, appears in Sturm--Liouville orthogonality and adjoint operators. Substitution changes variables using $\int f(g(x))g'(x)\dd x=\int f(u)\dd u$. Partial fractions are useful for inverse Laplace transforms and rational separable ODEs.
\begin{example}[Integral with a rational nonlinearity]
For $y'=y(2-y)$, on a non-equilibrium branch
\[\int\frac{\dd y}{y(2-y)}=\int\dd t,\qquad\frac12\log\abs{\frac{y}{2-y}}=t+C.\]
Rearranging yields $y(t)=2/(1+Ce^{-2t})$ for appropriate $C$. The equilibrium functions $y\equiv0$ and $y\equiv2$ should be checked separately rather than lost in division.
\end{example}
\section{Taylor series, convergence and approximation}
The expansion about $t_0$ is $y(t)=\sum_{k=0}^\infty y^{(k)}(t_0)(t-t_0)^k/k!$ when the function is analytic. A Taylor polynomial gives a local approximation, but a formal power series need not converge without additional hypotheses. If $y'=f(t,y)$, higher derivatives can be computed recursively, e.g. $y''=f_t+f_y f$. Taylor methods use these derivatives; the textbook treats the Taylor series method and the method of Frobenius separately (Chapters 3 and 8).
\section{Calculus of several variables and PDE notation}
For $u(x,t)$, $u_t$ and $u_{xx}$ denote partial derivatives. The gradient $\nabla u$ and Laplacian $\Delta u=\sum_i\partial^2u/\partial x_i^2$ extend to spatial dimensions. The product, chain and divergence identities help derive conservation laws. Boundary conditions may be \emph{Dirichlet} ($u$ prescribed), \emph{Neumann} (normal derivative prescribed), or \emph{Robin} (linear combination prescribed).
\chapter{Linear algebra and functional prerequisites}
\begin{bookground}Correspondence: Chapters 5, 6 and 9 (systems and matrix methods); Chapters 10--11 (orthogonality and eigenfunctions). Chapter 9 explicitly reviews matrices and vectors.\end{bookground}
\section{Vectors, matrices and eigenvalues}
A state-space ODE has $\mathbf x'=A\mathbf x+\mathbf b$. An eigenpair satisfies $A\mathbf v=\lambda\mathbf v$. If $A$ has independent eigenvectors $\mathbf v_i$, the homogeneous solution is $\sum_i c_i e^{\lambda_i t}\mathbf v_i$. The real part of $\lambda_i$ governs the exponential magnitude of each mode. Non-diagonalizable matrices may require generalized eigenvectors or the matrix exponential.
\begin{example}[Two-mode linear system]
For $A=\begin{psmallmatrix}-1&0\\0&-3\end{psmallmatrix}$ and $\mathbf x(0)=(2,-1)^T$, $\mathbf x(t)=(2e^{-t},-e^{-3t})^T$. Direct substitution yields $\mathbf x'=A\mathbf x$. Both eigenvalues are negative, so the origin attracts all solutions.
\end{example}
\section{Inner products and orthogonality}
The Euclidean inner product is $\mathbf u^T\mathbf v$ for real vectors. A weighted function inner product is
\[\langle f,g\rangle_w=\int_a^b f(x)g(x)w(x)\dd x,\qquad w(x)>0.\]
Orthogonality of distinct eigenfunctions under this inner product makes separation of variables and eigenfunction expansions practical. An \emph{orthonormal} set has pairwise inner products $\delta_{mn}$. Completeness and convergence require hypotheses that should not be inferred from orthogonality alone.
\section{Operator viewpoint}
A linear differential operator $L$ obeys $L(af+bg)=aL(f)+bL(g)$. The kernel $\ker L$ is the homogeneous solution space. To solve $Ly=f$, one seeks a particular solution plus a homogeneous correction satisfying the boundary data. The presence of a nontrivial null space may prevent existence for some forcings or destroy uniqueness; the Fredholm alternative in Chapter 11 makes this precise.
\chapter{Existence, uniqueness, and dependence on data}
\begin{bookground}Correspondence: Chapter 13, Sections 13.1--13.4. The statements below emphasize the standard local Lipschitz interpretation and do not substitute for the textbook's complete proofs.\end{bookground}
\section{Why a solution formula is not enough}
One must separate \emph{existence} (at least one solution), \emph{uniqueness} (at most one), \emph{continuation} (extension to a larger interval), and \emph{stability/continuous dependence} (behavior under perturbations). Consider $y'=\sqrt{\abs y}$ with $y(0)=0$: the zero solution exists, but delayed-start positive solutions also exist. Continuity of $f$ alone need not imply uniqueness.
\begin{theorem}[Common sufficient local criterion]
If $f(t,y)$ is continuous near $(t_0,y_0)$ and locally Lipschitz in $y$ there, then $y'=f(t,y)$, $y(t_0)=y_0$ has a unique local solution. Continuity of the partial derivative $f_y$ on a neighborhood is a sufficient (not necessary) condition for local Lipschitz continuity in $y$.
\end{theorem}
\section{Picard iteration}
Rewrite the IVP as the integral equation
\[y(t)=y_0+\int_{t_0}^{t}f(s,y(s))\dd s.\]
Starting at $y_0(t)=y_0$, define $y_{n+1}(t)=y_0+\int_{t_0}^{t}f(s,y_n(s))\dd s$. Under suitable hypotheses on a sufficiently short interval this sequence converges to the unique solution. For $y'=y$, $y(0)=1$, iterates give $1,\ 1+t,\ 1+t+t^2/2!,\ldots$, converging to $e^t$.
\section{Gronwall inequality and continuous dependence}
When an error envelope $e(t)$ satisfies $e(t)\le a+L\int_{t_0}^t e(s)\dd s$ for $t\ge t_0$, Gronwall's inequality gives $e(t)\le ae^{L(t-t_0)}$. This explains why perturbations remain controlled on fixed finite intervals under appropriate Lipschitz hypotheses, while sensitive systems can amplify errors over long times.
\begin{caution}The estimate is an upper bound, not a statement that every trajectory grows exponentially. A small residual may still correspond to a large solution error in a poorly conditioned BVP.\end{caution}
\chapter{Nondimensionalization and modeling discipline}
\begin{bookground}Correspondence: the modeling motivations of Chapters 3, 4, 5, 10 and 12. The systematic verification checklist and dimensionless scaling workflow are independent engineering extensions.\end{bookground}
\section{Dimensionless variables}
For $m y''+c y'+ky=0$, take $\tau=t\sqrt{k/m}$, $\zeta=c/(2\sqrt{mk})$, and a displacement scale $Y$. Writing $u(\tau)=y(t)/Y$ yields
\[u''+2\zeta u'+u=0,\]
where primes now mean differentiation in dimensionless time. This exposes the damping ratio and lets computations compare physically different oscillators with the same dimensionless response.
\section{Engineering assumptions register}
Document the dependent variable, units, parameter units, physically meaningful domain, boundary/initial data, external inputs, conserved quantities and constitutive laws. State why a lumped model is acceptable and what spatial heterogeneity it ignores. Distinguish mathematical stability, numerical stability, and physical robustness.
\section{Model validation hierarchy}
\begin{enumerate}
\item \textbf{Dimensional validity}: each additive term must share units.
\item \textbf{Analytic validity}: substitute a closed-form candidate back into the equation and data.
\item \textbf{Numerical verification}: refine time step or spatial mesh and inspect invariants.
\item \textbf{Model validation}: compare predicted observables with independent measurements and uncertainty.
\end{enumerate}
\chapter{Computational extension: executable mathematics}
\begin{modern}The following Python/SciPy workflow is new to this suite; the textbook's Chapter 1 includes Euler approximation and Chapter 3 develops numerical methods, but the specific library interface and QA checklist are not sourced from the textbook.\end{modern}
\section{Build a trustworthy IVP experiment}
\begin{lstlisting}
import numpy as np
from scipy.integrate import solve_ivp

def f(t, y):
    return -2.0 * y

t_eval = np.linspace(0, 3, 61)
sol = solve_ivp(f, (0, 3), [1.0], t_eval=t_eval,
                rtol=1e-9, atol=1e-12)
assert sol.success, sol.message
exact = np.exp(-2 * t_eval)
err = np.max(np.abs(sol.y[0] - exact))
assert err < 1e-7
print("maximum sampled error:", err)
\end{lstlisting}
The numerical solution should be checked against the initial value, an analytic benchmark when available, and a second tolerance. Use componentwise absolute tolerances for mixed-unit state variables. Store solver return status; an array of numbers is not evidence of successful convergence.
\section{Diagnostic questions}
\begin{itemize}
\item Does the direction field agree with the computed trajectories?
\item Are transformations valid over the full integration interval?
\item What happens to the numerical solution if the interval, tolerances, or initial data change?
\item Does the solution obey conservation, positivity, monotonicity, or endpoint constraints?
\item Is observed disagreement numerical error or model-form error?
\end{itemize}
\chapter{Mastery exercises and answers}
\section{Concept and calculation problems}
\begin{enumerate}
\item Classify $y''+t y'+y^3=0$ by order, linearity and homogeneity.
\item Solve $y'=3y$, $y(0)=2$, and state its maximal real interval.
\item Determine equilibria and local signs for $y'=y(y-4)$.
\item For $A=\operatorname{diag}(-2,1)$, describe stable and unstable eigendirections.
\item Verify that $y(t)=e^{-t}$ solves $y'+y=0$.
\item Identify why $y'=2\sqrt{\abs y}$ at $y(0)=0$ may violate uniqueness.
\item Derive the dimensionless parameters for $C T'=-h(T-T_a)$.
\end{enumerate}
\section{Answer key}
\begin{enumerate}
\item Second order; nonlinear because of $y^3$; the linear homogeneous/nonhomogeneous dichotomy is not applicable in its usual linear-operator sense.
\item $y=2e^{3t}$, interval $\R$.
\item Equilibria $0,4$; $y'>0$ outside $(0,4)$ and $y'<0$ within; $0$ locally attracting and $4$ repelling.
\item The $(-2)$ eigenvector direction decays and the $(1)$ direction grows; the origin is a saddle.
\item $y'=-e^{-t}$, hence $y'+y=0$ identically.
\item The right side is not locally Lipschitz in $y$ at zero; for example a delayed positive branch can be attached to the zero solution.
\item With $\tau=ht/C$, $\theta=(T-T_a)/\Delta T$ (constant $T_a$), $\theta'=-\theta$.
\end{enumerate}

\appendix
\chapter{References and reproducibility}
\begin{enumerate}
\item R.~K. Nagle, E.~B. Saff, and A.~D. Snider, \emph{Fundamentals of Differential Equations and Boundary Value Problems}, 7th ed., Pearson, 2018, ISBN 978-0-321-97710-6. Primary source for the \emph{topic map and terminology}; cited locally by chapter and section.
\item SciPy project, \emph{solve\_ivp} reference, \url{https://docs.scipy.org/doc/scipy/reference/generated/scipy.integrate.solve_ivp.html} (modern extension).
\item SciPy project, \emph{solve\_bvp} reference, \url{https://docs.scipy.org/doc/scipy/reference/generated/scipy.integrate.solve_bvp.html} (modern extension).
\item SciPy project, \emph{eigsh} reference, \url{https://docs.scipy.org/doc/scipy/reference/generated/scipy.sparse.linalg.eigsh.html} (modern extension).
\item Python, NumPy, SciPy, and SymPy project documentation; computational checks require recording package versions and tolerances.
\end{enumerate}
\section*{Verification record template}
\begin{tabularx}{\textwidth}{>{\bfseries}p{0.30\textwidth}X}\toprule
Problem identifier & \rule{0pt}{15pt}\ \\ \midrule
ODE/PDE and domain & \\[8pt]
Initial/boundary conditions & \\[8pt]
Units and parameter values & \\[8pt]
Analytic residual & \\ [8pt]
Grid/step refinement & \\ [8pt]
Independent solver comparison & \\ [8pt]
Version and reproducibility notes & \\ [8pt]\bottomrule
\end{tabularx}
\end{document}
